\chapter*{Abstract}
\addcontentsline{toc}{chapter}{Abstract}

Coronary Artery Disease (CAD), characterized by atherosclerotic luminal narrowing (stenosis) within epicardial coronary vessels, remains the foremost cause of global cardiovascular mortality. X-ray Coronary Angiography (XCA) represents the primary gold-standard diagnostic imaging modality for direct morphological assessment and percutaneous coronary intervention (PCI) guidance. However, routine visual assessment of coronary angiograms by interventional cardiologists exhibits significant inter-observer variability (reported between 15\% and 30\%), vulnerability to optical foreshortening, vessel overlap, low signal-to-noise fluoroscopic contrast, and lack of immediate 3D volumetric depth.

To address these diagnostic and interpretive challenges, this major project presents the design, empirical development, and technical integration of the \textbf{Coronary AI Workstation} --- an end-to-end intelligent medical software platform combining multi-stage deep learning pipelines with synchronized 3D interactive cardiovascular visual computing. The system integrates three specialized neural network architectures fine-tuned on the multi-center ARCADE (Automated Regional Coronary Angiography and Diagnostic Evaluation) challenge benchmark: (1) a \textbf{U-Net++} architecture with a pre-trained \textbf{ResNet34} backbone for full coronary arterial tree segmentation, attaining a validation Dice Similarity Coefficient (DSC) of 0.787 and Intersection over Union (IoU) of 0.655; (2) an \textbf{EfficientNet-B0} deep convolutional neural network for focal stenosis classification across candidate vessel regions of interest (ROIs), achieving a validation ROC-AUC of 0.995, accuracy of 96.80\%, and F1-score of 96.83\%; and (3) a dedicated \textbf{U-Net (ResNet34)} semantic segmentation network producing pixel-level stenosis boundary contours, evaluated across multiple decision thresholds with an optimal mean validation Dice score of 0.4926 and IoU of 0.3598 at threshold $\tau = 0.30$.

Model interpretability is natively integrated using Gradient-weighted Class Activation Mapping (Grad-CAM) targeting the deepest convolutional feature layers ($\_conv\_head$), validating anatomical attention alignment. Beyond 2D segmentation overlays, the platform implements an 8-connectivity connected-component spatial clustering algorithm that identifies multiple simultaneous lesion regions, assigns unique identifiers ($L1, L2, \dots, L_n$), and maps them onto anatomical coronary branches (LAD, LCX, RCA). Mapped lesions are translated into a synchronized suite of interactive visualizations rendered in Three.js and WebGL: a \textbf{3D Virtual Heart} displaying pulsing multi-beacon lesion anchors across coronary branches, an interactive \textbf{3D Coronary Artery Interior} featuring an anatomically scaled 180$^\circ$ longitudinal vessel cutaway with camera fly-through scrubbers, a multi-location \textbf{Cross-Section Inspector} providing lumen caliber comparisons (Minimal Lumen Area, proximal/distal references), an illustrative \textbf{Hemodynamic Particle Simulator} visualizing Bernoulli flow acceleration, an \textbf{AI Analytics Dashboard}, and automated \textbf{Clinical Research Case Reports}. The entire full-stack ecosystem connects a Python Flask REST service with a React/Vite front-end, fulfilling stringent data provenance, zero-lesion patent vessel fallback, and FDA Investigational Use Only (IUO) research governance.

\vspace{0.8cm}
\noindent
\textbf{Keywords:} Coronary Artery Disease, X-ray Coronary Angiography, Deep Learning, U-Net, U-Net++, EfficientNet-B0, Semantic Segmentation, Grad-CAM, 3D Cardiovascular Visualization, Three.js, Quantitative Coronary Analysis, Multi-Lesion Detection.
